% GENERATED FILE. Edit canonical lessons, then run npm run generate:books.
% canonical-source-hash: fnv1a64:a2ba64cc4766d7cb
% canonical-lessons: RU-C196-krysa, RU-C196-pauk, RU-C196-komar, RU-C196-mukha, RU-C196-zhuk

\chapter{Rat, Spider, Mosquito, Fly, Beetle}
\label{ch:ru-rat-spider-mosquito-fly-beetle}
\hlchaptermodality{196}

\begin{chapteropening}
\textbf{By the end of this chapter:} \emph{I can say a rat, a spider, a mosquito, a fly and a beetle.}

Complete the last lesson of chapter 196: I can say a rat, a spider, a mosquito, a fly and a beetle.
\end{chapteropening}

\section[\texorpdfstring{\ru{крыса} (krýsa)}{крыса (krýsa)}]{\ru{крыса} (krýsa) — a rat}

\label{lesson:RU-C196-krysa}



\begin{quote}
Before the new one: say the Russian for a rose, then the Russian for an oak.
\end{quote}

\subsection*{You'll want to know: \ru{крыса}}
\textbf{\ru{крыса}} (\emph{krýsa}) — \textquotedblleft{}a rat\textquotedblright{}.

\begin{grammarlens}[title={What you've built}]
One more word for this chapter.
\end{grammarlens}

\subsection*{Your turn}
\begin{itemize}
\raggedright
  \item \emph{Say it:} \emph{krýsa}
  \item \emph{Say it:} \emph{krýsa}, once more
  \item \emph{Say it:} say \emph{krýsa}
  \item \emph{Recall:} say the Russian for a rose, then the Russian for an oak, then say \emph{krýsa} again
\end{itemize}

\begin{tcolorbox}[breakable,colback=teal!4,colframe=teal!35!black,title={Before you move on}]
What does \ru{крыса} mean? (\textquotedblleft{}A rat\textquotedblright{}.) Say it once more.
\end{tcolorbox}
\section[\texorpdfstring{\ru{паук} (paúk)}{паук (paúk)}]{\ru{паук} (paúk) — a spider}

\label{lesson:RU-C196-pauk}



\begin{quote}
Before the new one: say the Russian for an oak, then the Russian for a rat.
\end{quote}

\subsection*{You'll want to know: \ru{паук}}
\textbf{\ru{паук}} (\emph{paúk}) — \textquotedblleft{}a spider\textquotedblright{}.

\begin{grammarlens}[title={What you've built}]
One more word for this chapter.
\end{grammarlens}

\subsection*{Your turn}
\begin{itemize}
\raggedright
  \item \emph{Say it:} \emph{paúk}
  \item \emph{Say it:} \emph{paúk}, once more
  \item \emph{Say it:} say \emph{paúk}
  \item \emph{Recall:} say the Russian for an oak, then the Russian for a rat, then say \emph{paúk} again
\end{itemize}

\begin{tcolorbox}[breakable,colback=teal!4,colframe=teal!35!black,title={Before you move on}]
What does \ru{паук} mean? (\textquotedblleft{}A spider\textquotedblright{}.) Say it once more.
\end{tcolorbox}
\section[\texorpdfstring{\ru{комар} (komár)}{комар (komár)}]{\ru{комар} (komár) — a mosquito}

\label{lesson:RU-C196-komar}



\begin{quote}
Before the new one: say the Russian for a rat, then the Russian for a spider.
\end{quote}

\subsection*{You'll want to know: \ru{комар}}
\textbf{\ru{комар}} (\emph{komár}) — \textquotedblleft{}a mosquito\textquotedblright{}.

\begin{grammarlens}[title={What you've built}]
One more word for this chapter.
\end{grammarlens}

\subsection*{Your turn}
\begin{itemize}
\raggedright
  \item \emph{Say it:} \emph{komár}
  \item \emph{Say it:} \emph{komár}, once more
  \item \emph{Say it:} say \emph{komár}
  \item \emph{Recall:} say the Russian for a rat, then the Russian for a spider, then say \emph{komár} again
\end{itemize}

\begin{tcolorbox}[breakable,colback=teal!4,colframe=teal!35!black,title={Before you move on}]
What does \ru{комар} mean? (\textquotedblleft{}A mosquito\textquotedblright{}.) Say it once more.
\end{tcolorbox}
\section[\texorpdfstring{\ru{муха} (múkha)}{муха (múkha)}]{\ru{муха} (múkha) — a fly}

\label{lesson:RU-C196-mukha}



\begin{quote}
Before the new one: say the Russian for a spider, then the Russian for a mosquito.
\end{quote}

\subsection*{You'll want to know: \ru{муха}}
\textbf{\ru{муха}} (\emph{múkha}) — \textquotedblleft{}a fly\textquotedblright{}.

\begin{grammarlens}[title={What you've built}]
One more word for this chapter.
\end{grammarlens}

\subsection*{Your turn}
\begin{itemize}
\raggedright
  \item \emph{Say it:} \emph{múkha}
  \item \emph{Say it:} \emph{múkha}, once more
  \item \emph{Say it:} say \emph{múkha}
  \item \emph{Recall:} say the Russian for a spider, then the Russian for a mosquito, then say \emph{múkha} again
\end{itemize}

\begin{tcolorbox}[breakable,colback=teal!4,colframe=teal!35!black,title={Before you move on}]
What does \ru{муха} mean? (\textquotedblleft{}A fly\textquotedblright{}.) Say it once more.
\end{tcolorbox}
\section[\texorpdfstring{\ru{жук} (zhuk)}{жук (zhuk)}]{\ru{жук} (zhuk) — a beetle}

\label{lesson:RU-C196-zhuk}



\begin{quote}
Before the new one: say the Russian for a mosquito, then the Russian for a fly.
\end{quote}

\subsection*{You'll want to know: \ru{жук}}
\textbf{\ru{жук}} (\emph{zhuk}) — \textquotedblleft{}a beetle\textquotedblright{}.

\begin{grammarlens}[title={What you've built}]
One more word for this chapter.
\end{grammarlens}

\subsection*{Your turn}
\begin{itemize}
\raggedright
  \item \emph{Say it:} \emph{zhuk}
  \item \emph{Say it:} \emph{zhuk}, once more
  \item \emph{Say it:} say \emph{zhuk}
  \item \emph{Recall:} say the Russian for a mosquito, then the Russian for a fly, then say \emph{zhuk} again
\end{itemize}

\begin{tcolorbox}[breakable,colback=teal!4,colframe=teal!35!black,title={Before you move on}]
What does \ru{жук} mean? (\textquotedblleft{}A beetle\textquotedblright{}.) Say it once more.
\end{tcolorbox}
